% ==============================================================================
% Fichier : front/00_page_titre.tex
% Titre   : Page de titre - Forensics et Analyse de Vulnérabilités
% Auteur  : MINKA MI NGUIDJOÏ Thierry Emmanuel
% Version : 1.0 -- novembre 2024
% ==============================================================================

\begin{titlepage}
\thispagestyle{empty}

% --- Fond décoratif : bandeaux bleu marine et filets or ---
\begin{tikzpicture}[remember picture, overlay]
    \fill[CouleurPrimaire]
        (current page.north west) rectangle
        ([yshift=-3.6cm] current page.north east);
    \fill[CouleurAccent]
        ([yshift=-3.6cm] current page.north west) rectangle
        ([yshift=-4.0cm] current page.north east);
    \fill[CouleurPrimaire]
        (current page.south west) rectangle
        ([yshift=2.5cm] current page.south east);
    \fill[CouleurAccent]
        ([yshift=2.5cm] current page.south west) rectangle
        ([yshift=2.9cm] current page.south east);
    \node[align=center, text=white] at ([yshift=-1.8cm] current page.north) {%
        {\fontsize{11}{13}\selectfont
        \textsc{République du Cameroun - Paix, Travail, Patrie}}\\[0.3ex]
        {\color{CouleurAccent!60!white}\rule{0.4\paperwidth}{0.4pt}}\\[0.6ex]
        {\fontsize{10}{12}\selectfont\textbf{\institutionprimaire}}\\[0.2ex]
        {\fontsize{8}{10}\selectfont \tutelle}\\[0.2ex]
        {\fontsize{9}{11}\selectfont \unitepedago}};
    \node[fill=white, rounded corners=4pt, inner sep=2pt]
        at ([xshift=1.75cm, yshift=-1.8cm] current page.north west)
        {\includegraphics[height=2.6cm]{figures/logos/enspy_logo.jpg}};
    \node[fill=white, rounded corners=4pt, inner sep=2pt]
        at ([xshift=-1.75cm, yshift=-1.8cm] current page.north east)
        {\includegraphics[height=2.6cm]{figures/logos/unesco_logo.jpeg}};
    \node[align=center] at ([yshift=1.25cm] current page.south) {%
        {\color{white}\bfseries\large Année académique \anneescolaire}\\[0.3ex]
        {\color{CouleurAccent!70!white}\small
        Version 1.0 | Document pédagogique de l'UE \codeUE}};
\end{tikzpicture}

\vspace*{0.6cm}

\begin{center}
    {\color{CouleurPrimaire}\rule{0.85\textwidth}{1.2pt}}\\[1.2ex]
    {\fontsize{26}{32}\selectfont\bfseries\color{CouleurPrimaire}
        Forensics et Analyse\\[0.4ex]
        de Vulnérabilités}\\[0.8ex]
    {\color{CouleurAccent}\rule{0.85\textwidth}{2pt}}\\[0.5ex]
    {\color{CouleurPrimaire}\rule{0.85\textwidth}{0.5pt}}
\end{center}

\vspace{0.8cm}

\begin{center}
    {\fontsize{13}{16}\selectfont\itshape\color{CouleurSecondaire}
        De la trace à la faille~:\\
        enquêter, prouver, corriger}
\end{center}

\vspace{0.5cm}

\begin{center}
    \begin{tikzpicture}
        \node[draw=CouleurPrimaire, fill=CouleurPrimaire, rounded corners=6pt,
              inner xsep=14pt, inner ysep=6pt]{%
            \color{white}\bfseries\large \sigleCours{} - UE~\codeUE
        };
    \end{tikzpicture}
\end{center}

\vspace{1.0cm}

% --- Schéma de la discipline : de la préservation au rapport, en passant par
%     la vulnérabilité et sa correction ---
\begin{center}
\begin{tikzpicture}[
    etape/.style = {draw=CouleurBordDef, fill=CouleurDef, rounded corners=3pt,
        minimum width=2.5cm, minimum height=0.8cm, align=center,
        font=\small\bfseries\color{CouleurPrimaire}, inner sep=2pt},
    arrow/.style = {-{Stealth[length=4pt]}, thick, color=CouleurSecondaire}
]
\node[etape] (n1) at (0, 0)      {Préserver};
\node[etape] (n2) at (3.2, 0)    {Acquérir};
\node[etape] (n3) at (6.4, 0)    {Analyser};
\node[etape] (n4) at (6.4, -1.4) {Vulnérabilité};
\node[etape] (n5) at (3.2, -1.4) {Correction};
\node[etape] (n6) at (0, -1.4)   {Rapport};
\draw[arrow] (n1) -- (n2); \draw[arrow] (n2) -- (n3); \draw[arrow] (n3) -- (n4);
\draw[arrow] (n4) -- (n5); \draw[arrow] (n5) -- (n6); \draw[arrow] (n6) -- (n1);
\node[draw=CouleurPrimaire, fill=CouleurPrimaire, rounded corners=4pt,
      align=center, font=\footnotesize\bfseries\color{white}, inner xsep=8pt]
      at (3.2, -2.8)
      {Chaque conclusion~: \OUI{} / \NON{} / \NEK{}, la constatation séparée de l'analyse};
\end{tikzpicture}
\end{center}

\vspace{0.6cm}

\begin{center}
    {\large\bfseries\color{CouleurPrimaire}
        \auteurcours}\\[0.5ex]
    {\normalsize\color{CouleurSecondaire}
       Enseignant Chercheur | \laboratoire{} | ENSPY}\\[0.2ex]
    {\small\color{CouleurBordInfo}
        \institutionprimaire}\\[0.2ex]
    {\small\color{CouleurBordInfo}
        \href{mailto:maletyon@proton.me}{maletyon@proton.me}}
\end{center}

\end{titlepage}
